\documentclass{article}
\usepackage[a4paper]{geometry}
\usepackage{fancyhdr}
\pagestyle{fancy}
\lhead{CRISPR}
\rhead{Mai 2026}

\begin{document}
\section{CRISPR/Cas9} 
Mit \emph{CRISPR}/\emph{Cas9} können Gene bearbeitet werden.

In das Enzym Cas9 wird dabei ein kleiner RNA-Strang, welcher an einer bestimmtem Stelle der DNA binden kann, eingeführt. Mithilfe diesem kann das Enzym die entsprechende Stelle der DNA finden, den Doppelstrang teilen und, so wie das Enzym es macht, die DNA an der Stelle schneiden.

Ist zeitgleich ein anderer DNA-Doppelstrang-Abschnitt in der Nähe, so kann der Körper im Versuch die eigentliche DNA zu reparieren, diesen Abschnitt hinzufügen.
\end{document}